\documentclass[12pt,a4paper]{article}
\usepackage[utf8]{inputenc}
\usepackage[T1]{fontenc}
\usepackage[brazil]{babel}
\usepackage{lmodern}
\usepackage{geometry}
\geometry{left=3cm,right=2cm,top=3cm,bottom=2cm}
\usepackage{xcolor}
\usepackage{listings}
\usepackage{microtype}
\usepackage{enumitem}
\usepackage[hidelinks]{hyperref}
\usepackage{array}
\setlength{\parindent}{1.25cm}
\setlength{\parskip}{0.25em}
\setlist{itemsep=0.2em,topsep=0.3em}
\definecolor{codebg}{RGB}{247,247,247}
\definecolor{codecomment}{RGB}{35,110,65}
\definecolor{codekeyword}{RGB}{35,65,145}
\lstdefinelanguage{Pifari}{
  morekeywords={true,false,null,None,Int,Float,Bool,Text,List,Map,Set,if,else,elseif,for,while,break,continue,return,import,const,await,async,agent,match},
  sensitive=true,
  morecomment=[s]{@}{@},
  morestring=[b]",
  morestring=[b]',
  alsoletter={|:}
}
\lstset{
  language=Pifari,
  basicstyle=\ttfamily\small,
  keywordstyle=\color{codekeyword}\bfseries,
  commentstyle=\color{codecomment}\itshape,
  stringstyle=\color{red!60!black},
  backgroundcolor=\color{codebg},
  frame=single,
  rulecolor=\color{black!20},
  breaklines=true,
  showstringspaces=false,
  keepspaces=true,
  columns=fullflexible,
  inputencoding=utf8,
  extendedchars=true
}
\title{\textbf{Definição da Linguagem de Programação Pifari}\\
\large Especificação consolidada -- versão de trabalho}
\author{Documento de projeto}
\date{Outubro de 2026}
\begin{document}
\maketitle
\begin{abstract}
Este documento consolida a proposta da linguagem de programação Pifari, de extensão \texttt{.pif}, com foco em sintaxe concisa, controle de fluxo orientado a valores, composição de operações, programação concorrente e assíncrona e interoperabilidade com C, C++, .NET e Python. A especificação incorpora a ideia de iniciar condicionais diretamente por expressões booleanas e iterações por valores iteráveis, sem exigir sempre as palavras-chave convencionais \texttt{if} e \texttt{for}. O texto define regras lexicais, sintáticas e semânticas propostas, exemplos, diagnósticos e pontos que ainda dependem de prototipação. A Pifari permanece uma linguagem em desenvolvimento; este documento especifica o comportamento pretendido, não afirma a existência de uma implementação completa.
\end{abstract}
\tableofcontents
\newpage

\section{Identidade, propósito e princípios}
\subsection{Identidade}
O nome da linguagem é \textbf{Pifari}, sua extensão é \texttt{.pif} e sua inspiração nominal vem do pífano e da tradição musical nordestina. A metáfora do projeto é a coordenação de diferentes tarefas que executam partes independentes de uma composição.

\subsection{Propósito}
A Pifari é concebida como uma linguagem interpretada, concorrente, assíncrona e de propósito geral. Seu objetivo é simplificar a coordenação de tarefas, a composição de operações e a criação de sistemas com agentes, mantendo interoperabilidade nativa com ecossistemas existentes. O runtime deverá cuidar do escalonamento de tarefas e do uso de um conjunto de threads, sem confundir concorrência com paralelismo: tarefas concorrentes podem progredir intercaladamente, enquanto tarefas paralelas executam efetivamente ao mesmo tempo quando houver recursos e independência suficientes.

\subsection{Princípios de projeto}
\begin{itemize}
\item expressões e valores devem ser capazes de iniciar construções de controle quando seu tipo e contexto forem compatíveis;
\item a sintaxe deve favorecer leitura da esquerda para a direita e composição com \texttt{|>};
\item a indentação define os blocos e deve ser consistente;
\item o sistema de tipos combina inferência, declarações explícitas e verificações em tempo de execução;
\item valores compartilhados, mensagens, tarefas assíncronas e agentes são mecanismos de primeira classe;
\item conversões entre tipos distintos devem ser explícitas, salvo conversões seguras formalmente definidas;
\item falhas de tarefas devem ser isoladas e diagnosticáveis;
\item interoperabilidade não deve esconder os custos e as regras de propriedade entre runtimes.
\end{itemize}

\section{Modelo de execução}
O programa-fonte é um arquivo UTF-8 com extensão \texttt{.pif}. A implementação de referência será um interpretador. O processamento conceitual compreende leitura do arquivo, análise léxica, análise sintática, análise semântica e de tipos, preparação de vínculos externos e execução pelo runtime. A implementação poderá armazenar uma representação intermediária, mas isso não altera a semântica definida pela linguagem.

O runtime gerenciará um pool de threads e um escalonador de tarefas. A existência de uma operação assíncrona não implica que ela seja executada em paralelo em todos os casos. Dependências de dados, efeitos colaterais, recursos compartilhados e limites do ambiente deverão ser considerados. A execução de código externo poderá ocorrer no mesmo processo ou em processo separado, conforme o mecanismo de interoperabilidade escolhido.

\section{Regras lexicais}
\subsection{Codificação e identificadores}
O código-fonte usa UTF-8. Identificadores começam por letra ou sublinhado e podem conter letras, números e sublinhados. A linguagem diferencia maiúsculas de minúsculas. Para favorecer interoperabilidade, recomenda-se utilizar identificadores ASCII sem acentos na primeira versão, embora textos e comentários possam conter caracteres Unicode.

Exemplos de identificadores válidos: \texttt{idade}, \texttt{\_total}, \texttt{nota2}, \texttt{calcular\_media}. Identificadores não podem começar por número nem conter espaços ou hífens.

\subsection{Literais}
Os literais previstos são:
\begin{itemize}
\item inteiros, como \texttt{20} e \texttt{-3};
\item reais, como \texttt{3.14} e \texttt{-0.5};
\item booleanos, \texttt{true} e \texttt{false};
\item textos delimitados por aspas simples ou duplas, além da forma nativa sem aspas em txt(...);
\item listas, mapas e conjuntos, conforme as regras de construção de coleções;
\item valor nulo, escrito \texttt{null}; na fronteira com Python, o valor correspondente é \texttt{None}.
\end{itemize}
A gramática deverá definir escapes de texto, notação científica e a forma exata dos literais de mapas e conjuntos. Aspas simples e duplas são aceitas para preservar familiaridade sintática com C e Python, sem implicar compatibilidade integral entre as linguagens. A forma nativa txt(conteúdo sem aspas) também é aceita; nesse formato, o conteúdo é tratado como texto literal até o parêntese de fechamento correspondente, conforme as regras de balanceamento e escape.

\subsection{Comentários}
Todo comentário começa com o delimitador \texttt{@} e termina no próximo \texttt{@}. A mesma regra vale para comentários de uma linha e multilinha: o comentário pode ocupar parte de uma linha ou atravessar várias linhas, mas sempre precisa ser fechado por \texttt{@}. Não são necessários aspas nem marcadores especiais para comentários multilinha. Comentários não são aninhados na versão inicial; a ausência do delimitador final deve produzir erro léxico.

\begin{lstlisting}
@ Comentario de linha @
idade = 20

@ Comentario que atravessa varias linhas.
Pode documentar uma decisao ou um trecho de codigo. @
nome = "Ana"
nome_sem_aspas =: txt(Ana)
\end{lstlisting}

\subsection{Indentação e blocos}
A indentação é obrigatória para delimitar blocos. O arquivo pode usar espaços ou tabulações, mas um mesmo bloco não pode misturar os dois estilos. A implementação deverá sinalizar indentação inconsistente em vez de tentar adivinhar a estrutura. Quebras de linha separam instruções simples, exceto quando a gramática indicar continuação de expressão. Chaves não são necessárias para delimitar os blocos ordinários da linguagem.

\section{Variáveis, declarações e escopo}
\subsection{Declaração dinâmica e tipada}
A atribuição simples com \texttt{=} cria ou atualiza uma variável dinâmica, cujo tipo pode mudar de acordo com as regras de execução. A forma \texttt{nome =: tipo(valor)} declara uma variável com tipo explícito, cuja identidade de tipo não pode mudar depois da declaração. O valor inicial deve ser compatível com o tipo declarado. Os construtores tipados são escritos em minúsculas, como \texttt{int(...)}, \texttt{float(...)}, \texttt{bool(...)} e \texttt{txt(...)}; os nomes dos tipos abstratos são documentados como Int, Float, Bool e Text.

A construção \texttt{txt(conteúdo)} aceita texto sem aspas. Quando o conteúdo estiver entre aspas simples ou duplas, essas aspas delimitam um literal convencional. A forma sem aspas trata o conteúdo como texto literal, não como nome de variável ou expressão; para converter uma variável em texto, deve-se usar uma operação de conversão explicitamente definida. Aspas simples e duplas também podem ser usadas em literais textuais fora de \texttt{txt(...)}.

\begin{lstlisting}
idade = 10
idade = "dez"       @ Permitido para variavel dinamica @

pontos =: int(20)
pontos = 21          @ Continua sendo inteiro se a atribuicao for compativel @
@ pontos = "vinte" Erro de tipo: a variavel tipada nao muda de tipo @

nome =: txt(Maria)
saudacao =: txt(Ola, mundo!)
nome_python =: txt("Ana")
nome_c =: txt('Ana')
\end{lstlisting}

O símbolo \texttt{:=} continua iniciando uma declaração de classe, conforme a Seção 10. O operador \texttt{=:} introduz declarações tipadas de variáveis e também a seção de membros de uma classe, conforme o contexto sintático. A gramática deverá tornar essa distinção determinística; a implementação não pode depender de uma interpretação arbitrária.

\subsection{Constantes e escopo}
A forma \texttt{const} poderá marcar uma referência que não recebe novas atribuições. Cada variável pertence ao escopo do bloco, função, agente ou módulo em que foi declarada. O uso de identificador inexistente deve gerar diagnóstico. Regras de sombreamento de nomes deverão ser explícitas; a proposta inicial proíbe duplicação no mesmo escopo.

\section{Sistema de tipos}
\subsection{Tipos fundamentais}
\begin{center}
\begin{tabular}{>{\ttfamily}p{3cm}p{10cm}}
\textnormal{\textbf{Tipo}} & \textnormal{\textbf{Significado}}\\
\hline
Int & Inteiro de precisão definida pela implementação, com interoperabilidade prevista com inteiros C compatíveis.\\
Float & Número real de ponto flutuante.\\
Bool & Valor lógico \texttt{true} ou \texttt{false}.\\
Text & Sequência de caracteres Unicode.\\
List[T] & Coleção ordenada de valores do tipo T.\\
Map[K,V] & Associação entre chaves K e valores V.\\
Set[T] & Coleção de valores únicos.\\
Null & Ausência intencional de valor, escrita \texttt{null}.\\
\end{tabular}
\end{center}

Também poderão existir tipos de função, tarefa, agente, objeto, módulo e tipos definidos pelo programador. Os tipos primitivos deverão oferecer mapeamentos explícitos e documentados para tipos C equivalentes, sem presumir que todas as plataformas usem a mesma largura numérica.

\subsection{Inferência, verificações e conversões}
A Pifari adota um sistema híbrido: o analisador infere tipos e detecta incompatibilidades sempre que possível, e o runtime verifica operações que dependam de valores dinâmicos. Uma variável dinâmica pode receber valores de tipos diferentes; uma variável declarada com tipo explícito não pode mudar de tipo.

Conversões entre tipos distintos devem ser explícitas, por meio de construtores ou funções de conversão, exceto conversões seguras que venham a ser especificadas formalmente. Não deve haver conversão implícita arbitrária de texto para número, nem de qualquer objeto para booleano. Uma condição exige um valor de tipo \texttt{Bool}. Um valor nulo não é equivalente a \texttt{false}; sua semântica é distinta.

\section{Expressões e operadores}
Os operadores aritméticos previstos são \texttt{+}, \texttt{-}, \texttt{*}, \texttt{/}, \texttt{//}, \texttt{\%} e \texttt{**}. Os relacionais são \texttt{==}, \texttt{!=}, \texttt{<}, \texttt{<=}, \texttt{>} e \texttt{>=}. Os operadores lógicos são \texttt{and}, \texttt{or} e \texttt{not}. O operador \texttt{=} representa atribuição. O operador \texttt{|>} representa composição e, em contextos iteráveis específicos, encadeamento de iteração.

A precedência exata será fixada na gramática formal. Como princípio, chamadas e agrupamentos têm precedência alta; operações aritméticas precedem comparações; comparações precedem operadores lógicos; e a composição por \texttt{|>} associa-se da esquerda para a direita, exceto nos contextos de iteração descritos adiante. Parênteses podem explicitar a ordem de avaliação.

\begin{lstlisting}
resultado_a = 2 + 3 * 4
resultado_b = (2 + 3) * 4
dobro = 5 ** 2
\end{lstlisting}

\section{Controle de fluxo orientado a valores}
\subsection{Princípio geral}
Uma expressão pode iniciar uma construção de controle quando seu tipo determina inequivocamente a intenção. Se a expressão resultar em \texttt{Bool}, ela pode iniciar um bloco condicional. Se resultar em um valor iterável, pode iniciar uma iteração por meio do pipeline. Portanto, não é obrigatório escrever \texttt{if} ou \texttt{for} em todas as construções. Essa regra não transforma identificadores como \texttt{test} em palavras reservadas: são expressões comuns interpretadas conforme o valor e o contexto.

\subsection{Condicional iniciado por booleano}
A forma básica é uma expressão booleana seguida de um bloco indentado. O bloco é executado quando a expressão resulta em \texttt{true}. A cláusula \texttt{elseif} introduz uma nova condição, avaliada apenas se as condições anteriores forem falsas. A cláusula \texttt{else} fornece o caminho alternativo quando nenhuma condição anterior for satisfeita.

\begin{lstlisting}
test = true

test
    print("ola")
elseif outra_condicao
    print("outra condicao verdadeira")
else
    print("nenhuma condicao satisfeita")
\end{lstlisting}

O exemplo ilustra a forma pretendida; \texttt{outra\_condicao} é um identificador ilustrativo e precisa estar definido como \texttt{Bool}. As condições são avaliadas em ordem, uma única vez cada, até que uma seja verdadeira. Apenas o bloco selecionado é executado. Um valor não booleano não pode ser usado silenciosamente como condição.

A forma explícita convencional \texttt{if condição} poderá ser aceita como alternativa de legibilidade, mas a forma direta por expressão é parte central da proposta. Parênteses não são obrigatórios ao redor da condição. A indentação determina o corpo e a associação de \texttt{elseif} e \texttt{else} ao condicional mais próximo que ainda aceite essas cláusulas.

\subsection{Condicionais como expressões}
Uma construção condicional pode produzir um valor quando todos os caminhos relevantes resultam em valores compatíveis. O operador \texttt{|>} não deve ser presumido como retorno apenas por aparecer no bloco alternativo: seu significado é definido pelo contexto sintático. A forma exata de escrever um resultado condicional será validada por protótipo para evitar confusão entre bloco de controle e composição de operações.

\begin{lstlisting}
nota = 8
nota >= 7
    print("Aprovado")
else
    print("Reprovado")
\end{lstlisting}

\subsection{Laços convencionais}
A Pifari pode oferecer \texttt{while} para repetição condicionada e uma forma explícita de \texttt{for} como alternativa legível. \texttt{break} interrompe a repetição atual e \texttt{continue} avança para a próxima iteração. Essas palavras não substituem a forma de iteração por pipeline; oferecem uma alternativa para algoritmos em que o fluxo explícito é mais claro.

\section{Iteração por valores iteráveis e operador \texttt{|>}}
\subsection{Iteração concisa}
Quando uma expressão produz um valor que implementa a capacidade \textit{Iterable}, ela pode iniciar uma iteração. Na forma concisa, o primeiro \texttt{|>} associa o elemento atual a um identificador e o próximo \texttt{|>} encadeia uma operação aplicada a cada elemento. Assim, a coleção, a variável de iteração e a ação podem ser lidas da esquerda para a direita.

\begin{lstlisting}
numeros = [1, 2, 3]

numeros
    |> numero
    |> print()
\end{lstlisting}

A semântica pretendida é equivalente a percorrer os elementos de \texttt{numeros}, vincular cada elemento a \texttt{numero} e executar \texttt{print(numero)} uma vez por elemento. A variável de iteração tem escopo limitado ao corpo da iteração. A ordem de visita preserva a ordem definida pelo iterável; coleções sem ordem garantida não passam a ter ordem garantida apenas por serem usadas com \texttt{|>}.

\subsection{Regras semânticas da iteração}
\begin{itemize}
\item a expressão inicial deve produzir um iterável;
\item o identificador depois do primeiro \texttt{|>} nomeia o elemento atual;
\item operações posteriores do pipeline são avaliadas para cada elemento, salvo se uma operação estiver explicitamente definida como agregadora;
\item a variável de iteração não pode ser lida fora de seu escopo;
\item erros em uma operação devem apontar o elemento e a etapa do pipeline quando essas informações estiverem disponíveis;
\item paralelizar a iteração é uma decisão do runtime, permitida apenas quando a semântica e os efeitos colaterais não forem violados.
\end{itemize}

\subsection{Composição funcional geral}
Fora do contexto de iteração concisa, \texttt{|>} encaminha o resultado da expressão à operação seguinte. As operações de coleção incluem transformações, filtragem e agregação. Lambdas usam \texttt{->}; esse operador é reservado para expressões lambda e não é um operador geral de retorno.

\begin{lstlisting}
numeros = [1, 2, 3, 4, 5]
dobrados = numeros |> map(x -> x * 2)
pares = numeros |> filter(x -> x % 2 == 0)
soma = numeros |> reduce((acum, x) -> acum + x, 0)
\end{lstlisting}

A gramática deverá distinguir a composição funcional comum da forma de iteração por elemento. Se uma sequência de tokens puder ser interpretada de ambas as maneiras, a especificação final deverá priorizar uma regra determinística ou exigir parênteses/uma forma explícita. O compilador não deve decidir com base em heurísticas não documentadas.

\section{Funções, lambdas e módulos}
Funções aceitam parâmetros e podem produzir valores. A forma final de declaração de função será definida na gramática. O operador \texttt{->} é usado somente para lambdas.

\begin{lstlisting}
dobrar = x -> x * 2
resultado = dobrar(5)

import matematica
\end{lstlisting}

O sistema de módulos utiliza \texttt{import} para importar módulos Pifari e módulos Python. A interoperabilidade Python deverá permitir tanto integração embutida no processo quanto comunicação com um processo Python separado. A escolha do modo de integração deve ser explícita ou seguir configuração documentada, pois os modelos têm custos e limites distintos.

\section{Classes, atributos e métodos}
A declaração de classe usa \texttt{:=} para iniciar a classe e \texttt{=:} para introduzir a seção de membros. Esses operadores possuem papéis diferentes e não são intercambiáveis. Atributos e métodos são privados por padrão; a forma de declarar membros públicos deverá ser especificada de maneira explícita.

\begin{lstlisting}
Pessoa :=
    =:
        nome =: txt(Ana)
        idade =: int(20)

        apresentar = () ->
            print(nome)
\end{lstlisting}

O trecho é esquemático e registra a distinção entre o início da classe e a seção de membros; a gramática definitiva ainda deve fixar a forma de declarar métodos, parâmetros, construtores e visibilidade. A resolução de nomes de membros, herança, interfaces e despacho dinâmico também deverá ser formalizada antes da implementação estável.

\section{Importação, bibliotecas nativas e interoperabilidade}
\subsection{\texttt{import}}
A palavra-chave \texttt{import} importa módulos Pifari e módulos Python. O sistema deverá definir resolução de nomes, caminhos de módulos, cache, inicialização e tratamento de falhas de importação.

\subsection{\texttt{\#include} e bibliotecas nativas}
A diretiva \texttt{\#include} destina-se à inclusão de declarações e interfaces nativas compatíveis, especialmente cabeçalhos C/C++. Ela não deve ser confundida com o sistema de módulos Python. Como C, C++ e C\# têm modelos de compilação e interoperabilidade distintos, o interpretador deverá encaminhar cada tipo de integração para o mecanismo correto, em vez de tratar todos como cabeçalhos equivalentes.

\subsection{\texttt{Link to}}
\texttt{Link to} seleciona ou vincula bibliotecas e símbolos externos definidos pelo sistema operacional ou pela Pifari. Pode receber nomes lógicos de bibliotecas ou caminhos de arquivos. A ligação de bibliotecas C deve ocorrer antes da execução do programa, de modo que símbolos ausentes ou incompatíveis sejam detectados antes de iniciar as tarefas que dependem deles.

\subsection{Fronteiras entre runtimes}
Valores entre Pifari, Python, C/C++ e .NET serão convertidos por padrão quando houver conversão segura e definida. Compartilhamento direto de memória entre runtimes deve ser explícito, pois envolve propriedade, ciclo de vida, representação e sincronização. Exceções, erros e recursos liberados por um runtime devem ser traduzidos ou encapsulados sem ocultar sua origem.

\section{Concorrência, assincronia e tarefas}
Concorrência e paralelismo são conceitos distintos. Uma operação assíncrona cria uma tarefa que precisa ser iniciada pelo mecanismo definido pela linguagem. \texttt{await} aguarda a conclusão de uma tarefa e só pode ser usado em contextos permitidos. O runtime administra o escalonamento e pode executar tarefas independentes em threads diferentes.

Falhas de uma tarefa devem ser isoladas daquela tarefa e comunicadas ao consumidor ou supervisor conforme a política de erros. O cancelamento é cooperativo e ocorre em pontos seguros. O runtime não deve interromper arbitrariamente uma operação nativa que não suporte cancelamento seguro.

\section{Agentes e comunicação}
Agentes são unidades de execução com estado próprio, tarefas e capacidade de trocar mensagens. Mensagens e memória compartilhada são mecanismos de primeira classe, com regras distintas. A comunicação por mensagens favorece isolamento; memória compartilhada pode ser usada quando a aplicação precisar de acesso coordenado a estado comum. O runtime deverá proteger automaticamente o acesso compartilhado conforme o modelo híbrido adotado, sem prometer ausência de bloqueios ou custo zero.

A criação, inicialização, encerramento, supervisão, política de falhas e endereçamento de agentes deverão ser formalizados em uma seção de runtime da implementação. Nenhum agente deve permanecer executando silenciosamente após o encerramento do programa.

\section{Memória e recursos}
A gestão de memória é automática por padrão. Controles explícitos de baixo nível poderão existir para casos que necessitem de interoperabilidade ou desempenho, sujeitos a regras de segurança e propriedade. Objetos compartilhados entre tarefas exigem sincronização definida pelo runtime. Recursos externos, como arquivos, sockets, referências nativas e objetos Python, devem ter regras claras de liberação.

O encerramento do programa aguarda as tarefas e os agentes que, segundo a política de ciclo de vida, devam concluir ou terminar. A política para tarefas destacadas, cancelamento no encerramento e falhas não tratadas deve ser determinística e documentada.

\section{Regras semânticas e diagnósticos}
O interpretador deverá distinguir erros léxicos, sintáticos, semânticos, de tipo e de execução. Mensagens devem indicar categoria, linha e coluna quando disponíveis, identificador ou expressão envolvida e uma explicação acionável.

\begin{itemize}
\item identificador não declarado: erro semântico;
\item expressão não booleana iniciando condicional: erro de tipo;
\item expressão não iterável iniciando iteração: erro de tipo;
\item uso de \texttt{elseif} ou \texttt{else} sem condicional compatível: erro sintático ou semântico;
\item mistura de tabulações e espaços na indentação do mesmo bloco: erro de indentação;
\item símbolo nativo não encontrado: erro de ligação anterior à execução;
\item conversão inválida: erro de conversão explícito;
\item falha de tarefa: erro associado à tarefa, propagado conforme a política de supervisão.
\end{itemize}

\section{Gramática preliminar}
A gramática abaixo é esquemática e registra os principais princípios. A gramática normativa deverá resolver a associação entre quebras de linha, indentação, blocos, pipeline funcional e pipeline iterável. Comentários de qualquer extensão seguem a regra única \texttt{@ ... @}.

\begin{lstlisting}
programa       := { instrucao }
instrucao      := atribuicao | expressao | condicional
                | iteracao | funcao | importacao
atribuicao     := IDENT "=" expressao
                | IDENT \"=:\" tipo \"(\" expressao \")\"
condicional    := expressao_bool BLOCO
                  { "elseif" expressao_bool BLOCO }
                  [ "else" BLOCO ]
iteracao       := expressao_iteravel "|>" IDENT
                  "|>" operacao
lambda         := parametros "->" expressao
importacao     := "import" IDENT
BLOCO          := NOVA_LINHA INDENT { instrucao } DEDENT
\end{lstlisting}

Essa gramática não é ainda uma EBNF completa: \texttt{expressao\_bool}, \texttt{expressao\_iteravel}, \texttt{operacao}, tipos, parâmetros e regras de indentação precisam ser definidos formalmente. A gramática final deverá refletir a distinção entre atribuição dinâmica com \texttt{=} e declaração tipada com \texttt{=:}, bem como a declaração de classe iniciada por \texttt{:=}.

\section{Exemplos integrados}
\subsection{Condição direta}
\begin{lstlisting}
ativo = true

ativo
    print("Sistema ativo")
else
    print("Sistema inativo")
\end{lstlisting}

\subsection{Iteração por pipeline}
\begin{lstlisting}
numeros = [1, 2, 3]

numeros
    |> numero
    |> print()
\end{lstlisting}

\subsection{Transformação e filtragem}
\begin{lstlisting}
numeros = [1, 2, 3, 4, 5]
dobrados = numeros |> map(x -> x * 2)
pares = numeros |> filter(x -> x % 2 == 0)
print(pares)
\end{lstlisting}

\subsection{Entrada, conversão e saída}
\begin{lstlisting}
texto = input("Digite sua idade: ")
idade = Int(texto)
print("Idade informada: ", idade)
\end{lstlisting}

Os exemplos são ilustrativos da proposta. A forma exata de chamadas, entrada/saída, visibilidade de membros, retorno de expressões condicionais e sintaxe de tarefas será confirmada por testes de parser e protótipos do runtime.

\section{Testes de conformidade iniciais}
A implementação deverá incluir testes que verifiquem ao menos:
\begin{enumerate}
\item booleano verdadeiro executa o primeiro bloco e booleano falso seleciona alternativa;
\item \texttt{elseif} avalia condições na ordem e \texttt{else} só é executado quando nenhuma condição é verdadeira;
\item valor numérico, texto ou nulo não é aceito como condição sem conversão explícita para \texttt{Bool};
\item iterável percorre os elementos conforme a ordem garantida pela coleção;
\item o primeiro estágio de iteração vincula o elemento ao identificador declarado;
\item a operação seguinte do pipeline é aplicada uma vez por elemento;
\item valor não iterável não pode iniciar iteração;
\item blocos com indentação inconsistente são rejeitados com diagnóstico;
\item variável tipada declarada com \texttt{=:} não muda de tipo, enquanto variável dinâmica segue as regras dinâmicas;
\item símbolos nativos ausentes são reportados antes da execução;
\item falha em uma tarefa não corrompe o estado de tarefas independentes;
\item cancelamento cooperativo e encerramento de agentes obedecem à política do runtime.
\end{enumerate}

\section{Decisões pendentes e critérios de estabilidade}
Esta especificação fixa a direção desejada, mas algumas decisões precisam de protótipos antes de serem consideradas estáveis:
\begin{itemize}
\item gramática completa para expressões que iniciam blocos e associação de \texttt{elseif}/\texttt{else};
\item distinção normativa entre pipeline funcional e pipeline de iteração;
\item sintaxe completa de classes, membros, visibilidade, métodos e construtores;
\item representação e largura dos tipos numéricos e seus mapeamentos para C;
\item modelo exato de importação Python embutida e por processo separado;
\item formato das declarações externas e resolução de símbolos de C/C++ e .NET;
\item sintaxe de criação, inicialização e supervisão de agentes;
\item regras finais de compartilhamento de memória, cancelamento e encerramento;
\item forma final de retorno de valores por blocos condicionais.
\end{itemize}

Uma regra será considerada estável quando possuir gramática não ambígua, semântica documentada, mensagens de erro definidas, testes de conformidade e implementação de referência consistente. A especificação deve distinguir claramente recursos normativos, recursos experimentais e ideias futuras.

\section{Conclusão}
A Pifari propõe uma linguagem em que valores e capacidades orientam a forma de controle: booleanos podem iniciar condicionais, iteráveis podem iniciar iterações por pipeline, e operações podem ser compostas da esquerda para a direita. A proposta combina essa sintaxe com inferência de tipos, verificações em tempo de execução, concorrência gerenciada, agentes, memória compartilhada protegida e interoperabilidade com ecossistemas C, .NET e Python. O próximo passo técnico é validar a gramática e as regras semânticas por meio de exemplos executáveis e testes automatizados, mantendo explícitas as decisões ainda experimentais.

\end{document}
